\documentclass[11pt,a4paper]{article}
\usepackage[utf8]{inputenc} % Not needed with XeLaTeX
\usepackage{fontspec}
\usepackage[margin=0.75in]{geometry}
\usepackage{xcolor}
\definecolor{darknavy}{RGB}{0,0,100}
\usepackage[colorlinks=true, urlcolor=darknavy, linkcolor=darknavy]{hyperref}
\usepackage{enumitem}
% \usepackage{fontawesome}
\usepackage{fontawesome5}
\usepackage{titlesec}
\usepackage{academicons}

% % --- Font: Source Sans Pro via Type1 (works with pdfLaTeX)
% \usepackage[default,scale=1.0]{sourcesanspro} % body + sans by default
% \renewcommand{\familydefault}{\sfdefault}      % make sans the default family
% \usepackage[scaled]{helvet}
\usepackage[default,scale=0.95]{opensans}

% (optional) math that pairs reasonably with sans body
\usepackage{newtxsf} % sans math; or use newtxtext/newtxmath for serif math

\titleformat{\section}{\Large\bfseries}{\thesection}{1em}{}[\titlerule]
\titlespacing{\section}{0pt}{12pt}{6pt}

\usepackage{comicneue}

\begin{document}

\begin{center}
  % {\Huge\textbf{Pavel Popov}}\\[0.5em]
  % {\comicneue\Huge\textbf{Pavel Popov}}\\[0.5em]
  {\fontspec{JetBrains Mono}\Huge Pavel Popov}\\[0.5em]
  {\fontspec{JetBrains Mono}\footnotesize +1 (470) 920-1876 \quad | \quad Atlanta, GA 30318} \quad | \quad {paavali.popov@gmail.com}\\
  {\small \href{https://github.com/paavalipopov}{Github: paavalipopov} \quad | \quad
  % \href{https://www.linkedin.com/in/pavel-popov-83b69a123/}{\faLinkedin\ Pavel Popov} \quad
  \href{https://www.linkedin.com/in/pavel-popov-83b69a123/}{LinkedIn: pavel-popov-83b69a123} \quad | \quad
  \href{https://scholar.google.com/citations?user=4alUD5UAAAAJ}{Google scholar: 4alUD5UAAAAJ}}
\end{center}

% Education first: the posting prefers MS/PhD candidates graduating Dec 2027 - May 2029.
\section*{Education}
\noindent\textbf{Georgia State University} \hfill Expected December 2027\\
PhD in Computer Science (MS in Computer Science awarded August 2026) \hfill GPA 3.89\\
{\small\textit{Coursework:} Advanced Machine Learning, Deep Learning, Advanced Data Mining, Big Data Programming}\\[0.3em]
\textbf{Moscow Institute of Physics and Technology} \hfill June 2020\\
MS in Applied Mathematics and Physics \hfill
% Spintronics,
GPA 3.92\\[0.3em]
\textbf{Moscow Institute of Physics and Technology} \hfill June 2018\\
BS in Applied Mathematics and Physics \hfill
% Solid State Physics,
GPA 3.81

\section*{Skills}
\begin{itemize}
[leftmargin=*,noitemsep]
\item \textbf{Machine Learning \& AI:} Deep learning (PyTorch, transformers), time-series modeling, regression, decision trees \& ensembles, clustering, model evaluation \& benchmarking, interpretable ML.
\item \textbf{LLMs \& Generative AI:} Working with AI agents; familiar with prompt engineering and LLMOps.
\item \textbf{Data \& Programming:} Python (Pandas, NumPy, SciPy, Scikit-learn, Matplotlib, Seaborn), SQL, NoSQL (mongoDB), Spark, Git, Docker, Linux/HPC (SLURM).
\item \textbf{Analytics \& Communication:} Excel, PowerPoint, data cleaning \& EDA, data visualization, technical writing, presenting to technical \& general audiences.
\end{itemize}

\section*{Open Source Projects}
\noindent\textbf{\href{https://pypi.org/project/ml4fmri/}{ml4fmri}: Python toolkit for benchmarking ML models on large-scale time-series data}\\
[-2em]
\begin{itemize}[leftmargin=*,noitemsep]
    \item Unified 13 ML models under a standardized training/evaluation API for reproducible comparison.
    \item Automated an end-to-end pipeline: data ingestion, preprocessing, training, and evaluation.
\end{itemize}

\noindent\textbf{\href{https://pypi.org/project/brainbow/}{brainbow} --- Data visualization tool for multi-dimensional datasets}\\
[-2em]
\begin{itemize}[leftmargin=*,noitemsep]
    \item Built a command-line tool for visualizing and exploring complex 3D/4D (neuroimaging) data.
\end{itemize}

\section*{Experience}
\noindent{\textbf{Center for Translational Research in Neuroimaging and Data Science}} \hfill January 2022 --- Now\\
Graduate Research Assistant \hfill \textit{Atlanta, USA}\\
[-2em]
\begin{itemize}[leftmargin=*,noitemsep]
    \item Designed and trained deep neural architectures (including transformers) for large-scale multimodal imaging \& time-series data (4D fMRI --- volumetric image/video data).
    \item Developed a robust deep learning classifier for medical diagnostics, achieving state-of-the-art performance with a 10x speedup over existing benchmarks.
    \item Scaled model training across multi-GPU HPC clusters using distributed data parallel (SLURM).
    \item Designed an interpretable AI model to simulate multivariate time-series data, enabling the identification of novel biomarkers and data-driven insights.
    \item Coordinated work of cross-functional teams to develop ML tools useful in clinical practice.
    \item Co-authored 4 peer-reviewed ML papers (e.g., \textit{NeuroImage}, IEEE MLSP) and presented findings to technical and general audiences at international conferences.
\end{itemize}

\noindent\textbf{Institute of Radioengineering and Electronics} \hfill May 2017 --- September 2020\\
Research Assistant \hfill \textit{Moscow, Russia}\\
[-2em]
\begin{itemize}[leftmargin=*,noitemsep]
    \item Built mathematical models and ran large-scale numerical simulations of high-frequency signal dynamics in novel hardware; co-authored 5 peer-reviewed papers (e.g., \textit{Physical Review Applied}).
\end{itemize}

\end{document}
